% =====================================================================
%  17.1  Disjoint Set Union (Union-Find)              sec:c28-dsu
%  Exports: def:c28-dsu  def:c28-find  prop:c28-alpha  prop:c28-cycle
%           ex:c28-merge  fig:c28-forest  lst:c28-find  lst:c28-union
%           lst:c28-comp   prob:c28-core  lst:c28-core-soln
% =====================================================================
\section{Disjoint Set Union: Parent, Rank, Path Compression}\label{sec:c28-dsu}

\begin{topicintro}
A \term{disjoint set union} (DSU, also \term{union-find}) maintains a
partition of \(n\) labelled items into \term{components} (sets) and supports
two operations: \code{find}, ``which set owns item \(v\)?'', and \code{union},
``glue the sets of \(u\) and \(v\) together.''  It is the data structure
behind Kruskal's MST, cycle detection, connected components, and every
``how many groups are left'' question -- and Chapters~6, 13, 14, 16 and 31
all called it without ever opening it.  This section owns the internals.

\textbf{Prerequisites.} \secref{c28-dsu} needs nothing but a rooted-tree
mental model; the DSU's own \code{parent} pointers are unrelated to any tree
\(T\) (they form their \emph{own} forest).  Arrays and recursion from
Chapter~1.

\textbf{After this section you can}
\begin{itemize}[nosep]
  \item implement \code{find} with path compression and \code{union} by size; \hfill\lv{2}
  \item state and justify the \(\bigO{\alpha(n)}\) amortised guarantee and
        know exactly when compression is legal; \hfill\lv{3}
  \item count components and reconnect a forest with the minimum new edges; \hfill\lv{3}
  \item detect a redundant edge (a cycle) purely by a \code{find} test. \hfill\lv{3}
\end{itemize}
\end{topicintro}

\begin{recognitionbox}[Pattern card -- ``glue sets together, then ask who owns what'']
\reckey{``minimum roads to connect all cities''; ``how many connected
components after these operations''; ``does adding this edge create a
cycle''; ``are \(u\) and \(v\) in the same group'' under a stream of
\emph{merges}; ``merge these two teams and keep track of the groups.''}
\recuse{
  \item only \emph{merging} (adding edges / joining sets) and asking about the
        resulting partition -- \secref{c28-dsu}, CSES~1666;
  \item cycle / ``is this edge redundant'' tests inside a growing graph --
        \propref{c28-cycle}, the Ch.~6 ``is it a tree?'' pre-look;
  \item a relative label per item appears -> upgrade to weighted DSU
        (\secref{c28-weighted}); undo / time appears -> rollback
        (\secref{c28-offline}).
}
\recnotuse{
  \item edges are \emph{deleted} online and connectivity is queried live --
        not monotone; DSU cannot un-merge unless you use the rollback
        variant offline (\(\secref{c28-offline}\));
  \item the graph structure per component (not just ``same set'') is needed
        -- maintain adjacency separately; the DSU only answers membership;
  \item a minimum spanning \emph{tree} is requested -- Kruskal uses DSU as a
        black box, but the cut-property reasoning is Ch.~31 (\chref{c17}).
}
One-line test: \emph{can the whole task be phrased as ``merge two groups, ask
which group owns a member,'' with merges never undone?  Then a plain DSU is
the structure.}
\end{recognitionbox}

% ---------------------------------------------------------------
\subsection{Intuition: A Forest of Parent Pointers}
% ---------------------------------------------------------------

Picture the \(n\) items as \(n\) single-node trees.  A \code{union} makes the
root of one tree a child of the root of another; the item at the very top
(whose \(\code{parent}[v]=v\)) is the set's \term{representative}.
\code{find(\(v\))} simply walks parent pointers upward until it reaches that
root.  Two items are in the same set iff their \code{find} results agree.
Nothing more exotic is happening -- the entire power comes from two tricks
that keep the trees shallow.

\begin{figure}[H]
\centering
\begin{tikzpicture}[diagram, scale=0.9, transform shape,
    dnode/.style={circle, draw=primaryblue, line width=0.9pt, fill=primaryblue!8,
                  minimum size=7mm, font=\small\bfseries, text=primaryblue!60!black},
    rnode/.style={circle, draw=accentorange, line width=1.1pt, fill=accentorange!16,
                  minimum size=7mm, font=\small\bfseries, text=accentorange!70!black}]
  % component A, root 1
  \node[rnode] (a1) at (0,0) {1};
  \node[dnode] (a2) at (-1.3,-1.3) {2};
  \node[dnode] (a4) at (-1.3,-2.6) {4};
  \node[dnode] (a3) at (1.3,-1.3) {3};
  \node[dnode] (a5) at (1.3,-2.6) {5};
  \draw[tedge] (a1)--(a2) (a2)--(a4) (a1)--(a3) (a3)--(a5);
  % component B, root 7
  \node[rnode] (b7) at (5,0) {7};
  \node[dnode] (b6) at (3.9,-1.3) {6};
  \node[dnode] (b8) at (6.1,-1.3) {8};
  \draw[tedge] (b7)--(b6) (b7)--(b8);
  \node[annot, font=\scriptsize, anchor=north] at (0,-3.2) {set with representative 1: \(\{1,2,3,4,5\}\)};
  \node[annot, font=\scriptsize, anchor=north] at (5.0,-2.0) {set with representative 7: \(\{6,7,8\}\)};
\end{tikzpicture}
\caption{The DSU's own forest after several unions -- not the tree \(T\), but
the partition \(\bigl\{\{1,2,3,4,5\},\{6,7,8\},\{9\},\{10\}\bigr\}\).  Orange
nodes are representatives (\(\code{parent}[v]=v\)); \code{find}(4) climbs
\(4\to2\to1\).  Shallow trees are the whole game: a deep chain makes
\code{find} linear, so \code{union} must not create deep chains.}
\label{fig:c28-forest}
\end{figure}

\begin{insightbox}[Two tricks, both about depth]
\textbf{Union by size/rank} attaches the \emph{smaller} tree under the
\emph{taller} root, bounding depth by \(\bigO{\log n}\) even with no
compression.  \textbf{Path compression} re-parents every node on a
\code{find} path straight onto the root, so the path is traversed the ``long
way'' only once.  Together they give an amortised \(\bigO{\alpha(n)}\); each
alone is weaker (union-by-size alone: \(\bigO{\log n}\); compression alone,
without union by rank: \(\bigO{\log n}\) amortised too, but not the near-linear
bound).
\end{insightbox}

% ---------------------------------------------------------------
\subsection{Definitions}
% ---------------------------------------------------------------

\begin{definition}{Disjoint set union}{c28-dsu}
A DSU on \(\{1,\dots,n\}\) is arrays \(\code{parent}[1..n]\) and
\(\code{sz}[1..n]\) with the invariant: following \(\code{parent}\) from any
\(v\) reaches a node \(r\) with \(\code{parent}[r]=r\) (the set's
\term{representative}); \(\code{sz}[r]\) stores the set's size and is
meaningful only when \(r\) is a root.  Initially \(\code{parent}[v]=v\),
\(\code{sz}[v]=1\) for all \(v\): \(n\) singleton sets.
\end{definition}

\begin{definition}{Find, with path compression}{c28-find}
\(\code{find}(v)\) returns the representative of \(v\)'s set.  With
\term{path compression}, every node visited on the way is re-parented
directly to the root before returning.
\end{definition}

% ---------------------------------------------------------------
\subsection{The Almost-Linear Guarantee}
% ---------------------------------------------------------------

\begin{proposition}{The \(\alpha(n)\) bound}{c28-alpha}
With union by size (or rank) and path compression, a sequence of \(m\)
operations on \(n\) items runs in \(\bigO{(n+m)\,\alpha(n)}\) total time,
where \(\alpha\) is the inverse Ackermann function.  For every conceivable
input, \(\alpha(n)\le 4\), so the cost is linear for all practical purposes.
\end{proposition}
\begin{proofbox}
Union by size keeps each tree's depth \(\le\log n\) (a node's set at least
doubles whenever it stops being a root).  The full amortised analysis
(Holm--de Lichtenberg--Thorup; Tarjan) charges a \code{find} against the
number of times its path's pointers get ``promoted'' by compression; that
recurrence inverts to \(\alpha\).  We use the black-box consequence: with both
tricks, each operation is effectively \(O(1)\).
\end{proofbox}

\begin{warningbox}[Compression is a one-way street]
Path compression rewrites intermediate \code{parent} pointers to jump to the
root.  It \emph{destroys} the information a rollback needs (which nodes you
touched on the path).  Therefore: \textbf{a rollback DSU must not use
compression} (\(\secref{c28-weighted}\), \(\secref{c28-offline}\)).  Whenever
the task never asks to undo, use compression freely -- that is the case here.
\end{warningbox}

\begin{proposition}{Redundant edge / cycle test}{c28-cycle}
When uniting the endpoints of an edge \((u,v)\): if \(\code{find}(u)=\code{find}(v)\)
\emph{before} the union, the edge closes a cycle inside one component and is
redundant; otherwise it merges two components and reduces the component count
by exactly \(1\).
\end{proposition}
\begin{proofbox}
\(\code{find}(u)=\code{find}(v)\) means a \(u\)--\(v\) path already exists in
this component, so adding \((u,v)\) creates a cycle.  If the roots differ, the
two sets become one, so the count drops by one.  A graph on \(n\) vertices is a
\emph{tree} iff exactly \(n-1\) edges pass the ``roots differ'' test and the
final count is \(1\) -- the DSU proof of ``connected + no cycle'' that
Chapter~6 (\chref{c25}) uses on LC~684/685 as an ``is it a tree?'' pre-look.
\end{proofbox}

% ---------------------------------------------------------------
\subsection{Algorithms and Implementation}
% ---------------------------------------------------------------

\begin{lstlisting}[caption={find with path compression (recursive; iterative form avoids deep chains)}, label={lst:c28-find}]
int find(int v) {
    if (parent[v] == v) return v;
    return parent[v] = find(parent[v]);  // compress on the way back up
}
// iterative, identical effect, safe when depth could reach n:
int find_it(int v) {
    int r = v;
    while (parent[r] != r) r = parent[r];        // climb to root
    while (parent[v] != v) { int p = parent[v];   // second pass:
        parent[v] = r; v = p; }                   // re-parent to root
    return r;
}
\end{lstlisting}

\begin{lstlisting}[caption={union by size (attach the smaller root under the larger)}, label={lst:c28-union}]
bool unite(int u, int v) {        // returns false if u,v were already together
    u = find(u); v = find(v);
    if (u == v) return false;     // same set -> redundant edge (prop:c28-cycle)
    if (sz[u] < sz[v]) { int t = u; u = v; v = t; }  // keep u the bigger root
    parent[v] = u;                // hang v's tree under u
    sz[u] += sz[v];               // u is the only root now; sz[v] stale
    return true;                  // two components merged -> count -= 1
}
\end{lstlisting}

\begin{lstlisting}[caption={component counting: start at n, decrement on every real merge}, label={lst:c28-comp}]
int components = n;
for (int i = 0; i < n; i++) { parent[i] = i; sz[i] = 1; }
for (auto [a, b] : edges)
    if (unite(a, b)) components--;   // only a genuine merge lowers the count
\end{lstlisting}

\begin{complexitybox}[Plain DSU]
\textbf{Per operation:} \(\bigO{\alpha(n)}\) amortised with both tricks.
\textbf{Build \(n\) singleton sets:} \(\bigO{n}\).  \textbf{\(m\) merges/finds:}
\(\bigO{(n+m)\,\alpha(n)}\) -- effectively linear.  \textbf{Memory:} \(\bigO{n}\).
\textbf{Recursion depth:} worst-case \(O(\log n)\) with union-by-size, so the
recursive \code{find} is safe here (unlike a DFS on a skewed tree).
\end{complexitybox}

% ---------------------------------------------------------------
\subsection{Worked Example: Merging the Edges of \(T\)}
% ---------------------------------------------------------------

\begin{example}{Union-by-size over the edges of the running tree}{c28-merge}
Process \(T\)'s nine edges in their input order
\((1\!-\!2),(1\!-\!3),(2\!-\!4),(2\!-\!5),(3\!-\!6),(3\!-\!7),(5\!-\!8),
(7\!-\!9),(9\!-\!10)\), starting from ten singletons.  \(\code{parent}[v]\)
below is written ``\(v\!\to\!p\)''; a root keeps \(\code{parent}[v]=v\).
\end{example}
\begin{tracebox}[parent / sz after each union on T]
\[
\renewcommand{\arraystretch}{1.2}
\begin{array}{c|l|l|l}
\text{step} & \text{unite} & \text{action} & \text{roots' sizes after} \\ \hline
0 & - & 10\text{ singletons} & \text{all } \code{sz}=1\\
1 & (1,2) & 2\to1,\ \code{sz}[1]=2 & 1\!:\!2\\
2 & (1,3) & 3\to1,\ \code{sz}[1]=3 & 1\!:\!3\\
3 & (2,4) & \code{find}(2)=1;\ 4\to1,\ \code{sz}[1]=4 & 1\!:\!4\\
4 & (2,5) & 5\to1,\ \code{sz}[1]=5 & 1\!:\!5\\
5 & (3,6) & \code{find}(3)=1;\ 6\to1,\ \code{sz}[1]=6 & 1\!:\!6\\
6 & (3,7) & 7\to1,\ \code{sz}[1]=7 & 1\!:\!7\\
7 & (5,8) & 8\to1,\ \code{sz}[1]=8 & 1\!:\!8\\
8 & (7,9) & 9\to1,\ \code{sz}[1]=9 & 1\!:\!9\\
9 & (9,10) & 10\to1,\ \code{sz}[1]=10 & 1\!:\!10
\end{array}
\]
Every \(\code{find}\) resolved to root~\(1\) by step~9, all nine edges merged
distinct sets (so \(\code{unite}\) returned \texttt{true} each time), the
component count went \(10\to1\), and \emph{no} edge triggered the cycle test
-- as required, since \(T\) is a tree with exactly \(n-1=9\) edges.  Had an
edge \((6,8)\) been added, \(\code{find}(6)=\code{find}(8)=1\) and
\(\code{unite}\) returns \texttt{false}: a detected cycle.
\end{tracebox}

% ---------------------------------------------------------------
\subsection{Problem Lab}
% ---------------------------------------------------------------

\solvedproblem{c28-core}{Building Roads}{CSES 1666}
  {https://cses.fi/problemset/task/1666}{CSES L3}
  {disjoint set union, component counting}

\begin{problemstatement}[CSES 1666 Building Roads]
Byteland has \(n\) cities and \(m\) existing roads (a road connects two
different cities; at most one road per pair).  Build the \emph{minimum} number
of new roads so that every city is reachable from every other, and print which
roads to build.

\textbf{Input.} \(n\) and \(m\) (\(1\le n\le 10^5\), \(1\le m\le 2\cdot 10^5\));
then \(m\) lines, each a road \(a\ b\).

\textbf{Output.} First line: the number \(k\) of new roads.  Then \(k\) lines,
each naming the two cities to join (any valid set).

\textbf{Example.}
\begin{verbatim}
Input:
4 2
1 2
3 4

Output:
1
2 3
\end{verbatim}
\end{problemstatement}

\begin{attemptgate}
\peek{``fewest roads to connect everything'' -- that is components minus one.}
\nudge{Run a DSU over the existing roads; each surviving component
contributes exactly one representative; chaining the representatives gives
the answer.}
\fullhint{Merge all \(m\) roads, collect one root per distinct
\(\code{find}\), print \(\text{count}-1\) and pair the first representative
with every other.}
\end{attemptgate}

\solstep{Understanding}
The existing roads partition the cities into connected components.  One new
road can join at most two components, so connecting \(c\) components needs at
least \(c-1\) roads -- and \(c-1\) suffice by chaining representatives.  The
task is exactly ``count components and emit a spanning set of them,'' a direct
\secref{c28-dsu} application.

\solstep{Key Observations}
\begin{observation}
A road whose endpoints are already in the same component (cycle test,
\propref{c28-cycle}) changes nothing -- the DSU simply ignores it.
\end{observation}
\begin{observation}
After all merges, the cities \(v\) with \(\code{find}(v)=v\) are the
\emph{distinct} component representatives; if there are \(c\) of them, the
answer is \(c-1\) roads.
\end{observation}
\begin{observation}
Pairing the \emph{first} representative with each other representative
\((r_1,r_i)\) connects all components in a star -- always valid, since joining
two different components never closes a cycle.
\end{observation}

\solstep{Think It Through}
\begin{thinkbox}
\deadend{Do a BFS/DFS per unvisited city to count components}
{Works, but then you must also record one representative per component and
re-scan; on \(n\le 10^5\) with \(m\le 2\cdot10^5\) it is fine but clunkier than
a DSU, which counts and picks representatives in the same pass.}
\aha{One DSU pass merges the roads and drops the count to \(c\); the roots of
the final forest are the representatives to chain.}
\end{thinkbox}

\solstep{Implementation}
\begin{lstlisting}[caption={CSES 1666 full solution (DSU + representative chain)}, label={lst:c28-core-soln}]
#include <bits/stdc++.h>
using namespace std;
const int MX = 100005;
int parent_[MX], sz_[MX];
int find(int v){ return parent_[v]==v ? v : parent_[v]=find(parent_[v]); }
bool unite(int a,int b){
    a=find(a); b=find(b);
    if(a==b) return false;
    if(sz_[a]<sz_[b]) swap(a,b);
    parent_[b]=a; sz_[a]+=sz_[b];
    return true;
}
int main(){
    int n,m; scanf("%d%d",&n,&m);
    for(int i=1;i<=n;i++){ parent_[i]=i; sz_[i]=1; }
    for(int i=0;i<m;i++){ int a,b; scanf("%d%d",&a,&b); unite(a,b); }
    vector<int> reps;
    for(int i=1;i<=n;i++) if(find(i)==i) reps.push_back(i);
    printf("%d\n",(int)reps.size()-1);
    for(size_t i=1;i<reps.size();i++) printf("%d %d\n",reps[0],reps[i]);
}
\end{lstlisting}

\solstep{Trace on the Official Sample}
\begin{tracebox}[CSES 1666: n=4, roads (1,2),(3,4)]
Start \(\code{parent}=(1,2,3,4)\), all \(\code{sz}=1\).
\(\text{unite}(1,2)\): roots \(1,2\); equal size, keep root \(1\); \(2\to1\),
\(\code{sz}[1]=2\).
\(\text{unite}(3,4)\): \(4\to3\), \(\code{sz}[3]=2\).
Representatives (roots): \(\{1,3\}\) -- \(c=2\).
Print \(c-1=1\), then the single new road \(1\ 3\) (a valid answer; the
sample's \texttt{2 3} is the same idea with representative \(2\) standing in
for the \(\{1,2\}\) component -- \code{find}(2)=1 here, so our output is
\texttt{1 3}, also accepted).
\end{tracebox}

\solstep{Complexity}
\begin{complexitybox}
\(n+m\le 3\cdot10^5\) operations, each \(\bigO{\alpha(n)}\) amortised: total
\(\bigO{(n+m)\alpha(n)}\).  Memory \(\bigO{n}\).  Well inside 1.00\,s.
\textbf{AC.}
\end{complexitybox}

\subsection{Common Pitfalls}

\begin{warningbox}[Trap 1: reusing stale \code{sz}]
After a union, \(\code{sz}[v]\) is only correct at the \emph{root}.  Reading
\(\code{sz}[x]\) for a non-root \(x\) gives a meaningless (frozen) number.
Always \(\code{sz}[\code{find}(x)]\).
\end{warningbox}

\begin{warningbox}[Trap 2: forgetting the cycle branch]
If you \code{union} unconditionally without checking \(\code{find}(u)=\code{find}(v)\),
you double-add \(\code{sz}\) and the component count drifts.  Route everything
through a \code{unite} that returns whether a real merge happened
(\coderef{c28-union}).
\end{warningbox}

\begin{warningbox}[Trap 3: 0- vs 1-indexing]
CSES numbers cities from \(1\).  Initialising the arrays as
\(\code{for }i=0..n-1\) while reading \(1..n\) leaves \(\code{parent}[n]\)
uninitialised.  Match the problem's indexing (this file uses \(1..n\)).
\end{warningbox}

% ---------------------------------------------------------------
\subsection{Practice Problems}
% ---------------------------------------------------------------

\begin{stuckbox}
\begin{itemize}[nosep, leftmargin=1.4em]
  \item Is every operation a merge (never an undo)?  If not, you need rollback
        (\(\secref{c28-offline}\)), not this DSU.
  \item Did you decrement the component count \emph{only} when \code{unite}
        returned \texttt{true}?
  \item Are you reading \code{sz} at a root, or at a raw node?
  \item Does your representative scan use \(\code{find}(i)=i\), not
        \(\code{parent}[i]=i\) (compression may have re-pointed the latter)?
\end{itemize}
\end{stuckbox}

\begin{practiceset}
  \practice{Building Teams}{CSES 1668}
    {https://cses.fi/problemset/task/1668}
    {CSES L3}
    {friend-enemy partition -- bipartite check per component; pair DSU with a
     2-colouring DFS, or use a weighted DSU (\(\secref{c28-weighted}\)).}
  \practice{Number of Islands II}{LC 305}
    {https://leetcode.com/problems/number-of-islands-ii/}
    {LC Hard L3}
    {cells turn to land one at a time; report the island count after each drop.  A
     new land cell starts its own component (count\,+\(1\)) and \code{unite}s with
     each land neighbour that merges (count\,\(-\,1\) per genuine merge) -- the
     monotone-additions pattern of \secref{c28-dsu} with \code{sz} driving the
     component tally.}
\end{practiceset}

\paragraph{Back-cites.}
Chapter~6 (\chref{c25}) -- the ``is it a tree?'' DSU pre-look on LC~547/684/685;
\propref{c28-cycle} is the internal proof it leans on.
Chapter~31 (\chref{c17}) -- Kruskal's MST calls \code{find}/\code{unite} as a
black box; the internals are here.

% ---------------------------------------------------------------
\begin{recallbox}[Section 17.1 in 10 seconds]
A DSU is a parent forest + a size stored at each root.  \code{find} climbs and
compresses; \code{union} glues the smaller root under the larger; both tricks
together give \(\bigO{\alpha(n)}\) amortised.  Component count starts at \(n\)
and drops by one per genuine merge; \(\code{find}(u)=\code{find}(v)\) before a
union means a redundant edge.  Compression is \emph{irreversible} -- never use
it if you must roll back.
\selftest{
  \item Write \code{find} with path compression from memory.
  \item Why must union-by-size attach small under large, and what depth does
        it guarantee?
  \item State the cycle test in one line of code.
  \item What breaks if you add path compression to a rollback DSU?
  \item On \(T\)'s edges above, what was \(\code{sz}[1]\) after step~5?  (6.)
}
\end{recallbox}
